%!TEX root = ../template.tex
%%%%%%%%%%%%%%%%%%%%%%%%%%%%%%%%%%%%%%%%%%%%%%%%%%%%%%%%%%%%%%%%%%%%
%% abstract-en.tex
%% NOVA thesis document file
%%
%% Abstract in English
%%%%%%%%%%%%%%%%%%%%%%%%%%%%%%%%%%%%%%%%%%%%%%%%%%%%%%%%%%%%%%%%%%%%

\typeout{NT FILE abstract-en.tex}%

Software and conceptual modelling are commonly included in software engineering and related computing programmes, but they often receive less emphasis than other topics in computing education. Students frequently experience modelling as either an essential aid for understanding and communication or as burdensome extra documentation. Existing research offers valuable insights but is dispersed across modelling traditions and study contexts, with limited integrative explanation of how challenges arise and how they relate to strategies, support, and course conditions.

This thesis investigates how software modelling is experienced, taught, and learned in higher education, treating modelling as a socio-technical activity. It addresses three research questions: (RQ1) what challenges students and instructors perceive and how they manifest across interpreting, creating, refining, and evaluating models; (RQ2) which strategies and pedagogical or tool-related support needs they identify; and (RQ3) how tools, course design, domains, and feedback practices shape those experiences.

The study used Socio-Technical Grounded Theory. Fourteen semi-structured interviews (four lecturers, ten students) were open-coded into 277 codes and 78 memos, then compared into 14 categories and four phenomena. Wave~2 theoretical sampling densified properties without adding a 15th overview category. The analysis claims local saturation of that category set and theoretical sufficiency for this sample, not global theoretical saturation.

The account shows a gated first modelling move distinct from symbol overload; difficulty judging ``good enough'' without a compile result; late and unequally accessible feedback, with rework depending on remaining time and locked grades; demand for checks and explanations rather than generated homework; and diagrams that survive only when someone owns them and someone reads them. Course load and tool friction condition whether examples and feedback loops actually occur. Implications for teaching and educational tools are derived from these patterns but were not trialled.

% Abstract keywords
\keywords{
  software modelling education \and
  conceptual modelling \and
  grounded theory \and
  socio-technical perspective \and
  modelling competence
}